% !TEX root = ../Main.tex
\chapter{混合泳与水上项目}

\section{个人混合泳}

\state{泳式顺序}{
    \textbf{个人混合泳}必须按照以下顺序比赛：
    \begin{center}
        \textbf{蝶泳 $\to$ 仰泳 $\to$ 蛙泳 $\to$ 自由泳}
    \end{center}
    每种泳式必须完成\textbf{赛程四分之一}的距离。
}

\section{混合泳接力}

\state{泳式顺序（与个人混合泳不同）}{
    \textbf{混合泳接力}必须按照以下顺序比赛：
    \begin{center}
        \textbf{仰泳 $\to$ 蛙泳 $\to$ 蝶泳 $\to$ 自由泳}
    \end{center}
    \textbf{注意}：混合泳接力\textbf{以仰泳开局}（因仰泳是水中出发，无法接棒于之前的泳式），其余顺序也与个人混合泳不同。
}

\section{混合泳通用规则}

\state{各泳式的合法性}{
    \begin{itemize}
        \item 在个人混合泳和混合泳接力项目中，\textbf{每一泳式都必须符合相关泳式的有关规定}（见第 4--7 章）；
        \item 在\textbf{仰泳转蛙泳}过程中，运动员必须\textbf{呈仰泳姿势触及池壁}。
    \end{itemize}
}

\section{两种顺序对比}

\begin{center}
\begin{tabular}{@{} >{\bfseries}l l @{}}
    \toprule
    项目 & 泳式顺序 \\
    \midrule
    个人混合泳 & 蝶 $\to$ 仰 $\to$ 蛙 $\to$ 自 \\
    混合泳接力 & 仰 $\to$ 蛙 $\to$ 蝶 $\to$ 自 \\
    \bottomrule
\end{tabular}
\end{center}

\section{游泳相关的水上项目分类}

\state{四大水上项目}{
    在\textbf{广义的游泳运动}下，包括以下 4 类项目：
    \begin{itemize}
        \item \textbf{游泳}（竞技游泳）
        \item \textbf{跳水}
        \item \textbf{花样游泳}
        \item \textbf{水球}
    \end{itemize}
}

\section{蛙泳动作要领（口诀）}

\state{四句十六字}{
    \begin{center}
        先\textbf{伸胳膊}再\textbf{蹬腿}，臂腿\textbf{伸直漂一会}。\\
        \textbf{低头伸臂}慢呼气，\textbf{划臂抬头}快吸气。
    \end{center}
    \begin{itemize}
        \item \textbf{动作顺序}：伸臂 $\to$ 蹬腿 $\to$ 伸直漂滑；
        \item \textbf{呼吸配合}：低头伸臂时\textbf{慢呼气}，划臂抬头时\textbf{快吸气}。
    \end{itemize}
}

\rmkb{
    \textbf{易考点}：
    \begin{itemize}
        \item \textbf{个人混合泳}：蝶 $\to$ 仰 $\to$ 蛙 $\to$ 自；\textbf{每种泳式占 1/4 赛程}；
        \item \textbf{混合泳接力}：仰 $\to$ 蛙 $\to$ 蝶 $\to$ 自（首棒为仰泳）；
        \item 仰 $\to$ 蛙 转身：必须\textbf{仰卧姿势触壁}；
        \item 水上项目 \textbf{4 大类}：游泳 / 跳水 / 花样游泳 / 水球；
        \item 蛙泳口诀：\textbf{先伸胳膊再蹬腿}；\textbf{低头伸臂慢呼气，划臂抬头快吸气}。
    \end{itemize}
}
